% paper/sections/11d_pure_fccd_dns.tex
%
% §11  純 FCCD DNS 構成（第1--4層）
%
% ═══════════════════════════════════════════════════════════════════════════════
\subsection{純 FCCD DNS 構成：シャープ界面極限への発展的閉包}
\label{sec:pure_fccd_dns}
% ═══════════════════════════════════════════════════════════════════════════════

本節では，標準経路で検証した分相 PPE，HFE，FCCD 面ジェット，
Balanced--Force 条件を一つの面基準点へ集約した場合に，
どの追加条件が必要になるかを整理する．
\textbf{純 FCCD シャープ界面 DNS}は本稿の標準経路そのものではなく，
高密度比，高粘性比，低 We が同時に成立するシャープ界面極限へ向けた
発展的閉包である．
したがって本節では期待性能ではなく，検証済み基盤量から
直接引き継げる条件と，本稿では未検証に留める範囲を分けて述べる．

\noindent\textbf{第1--4層の構成と適用範囲：}
\begin{center}
\small
\begin{tabular}{@{}llll@{}}
\hline
層 & 内容 & 主要技術 & 適用範囲 \\
\hline
第1層 & 界面追跡層 & Ridge--Eikonal + FCCD 輸送 & 標準界面追跡層として適用 \\
第2層 & 運動量移流層 & FCCD + HFE & 純 FCCD DNS 用 \\
第3層 & 粘性項分離 & 応力発散 + DC & 純 FCCD DNS 用 \\
第4層 & 分相 FCCD PPE & 分相 + GFM + DC $k{=}3$ & 純 FCCD DNS の中核 \\
\hline
\end{tabular}
\end{center}
第1層の軽量版は，PPE 分岐に依存しない標準界面追跡層として適用できる；
第1--4層の同時使用は本節の発展的構成を前提とする．

% -------------------------------------------------------------------------------
\subsubsection{設計命題：高次演算子の共通基準面位置}
\label{sec:pure_fccd_thesis}
% -------------------------------------------------------------------------------

純 FCCD DNS の中核命題：
\begin{quote}
\textbf{界面追跡・運動量移流・粘性・圧力投影・表面張力・HFE}
のすべての演算が，\textbf{共通の面基準点 $f$}で
FCCD 面ジェット $\FaceJet{\cdot}=(\cdot_f,\cdot'_f,\cdot''_f)$ を
基本量として高次一貫に定義される．
\end{quote}

この命題は §\ref{sec:bf_p1_location}（BF P-1 位置一致）を
全演算子に拡張したものに相当する．
純 FCCD DNS は (R1)-(R4)（§\ref{sec:fccd_def} の 4 役割定義）の面ジェット条件を
圧力・界面・移流・粘性の全演算子に
拡張した構成であり，単一の面ジェット構造で 4 経路の計量整合を
構造的に担保する．

% -------------------------------------------------------------------------------
\subsubsection{第1層：界面追跡層（Ridge--Eikonal + FCCD 輸送）}
\label{sec:pure_fccd_phase1}
% -------------------------------------------------------------------------------

界面は VOF 体積分率ではなく\textbf{高次輸送される連続場}として扱う：
\begin{itemize}[noitemsep]
  \item \textbf{トポロジー/計量分離}：
        Ridge--Eikonal ハイブリッド（§\ref{sec:ridge_eikonal}）が
        距離関数品質 ($\phi$ の符号距離性) とトポロジー変化
        （ピンチオフ，合体）を別層で扱う．
  \item \textbf{保存形輸送}：
        FCCD 保存形面フラックス（§\ref{sec:fccd_advection} の選択肢 B）で
        $\psi$ を進める．低次 VOF フラックスには置換しない．
  \item \textbf{幾何資産}：
        $\phi$，法線 $\hat{\bm n}_{lg}$，曲率 $\kappa_{lg}$ は
        $\psi$ から再生成され，
        表面張力・HFE・粘性法線-接線切替が共通で使用する．
\end{itemize}

界面幅 $\varepsilon$ は固定でなく\textbf{空間依存} $\varepsilon_\text{local}(\bx)$
（§\ref{sec:ridge_eikonal_nu_grid}）として
非一様格子に整合させる．Ridge--Eikonal が保証する符号距離性により
$\bnabla\cdot\hat{\bm n}_{lg}$ の高次評価が成立する．

% -------------------------------------------------------------------------------
\subsubsection{第2層：運動量移流層（FCCD + HFE）}
\label{sec:pure_fccd_phase2}
% -------------------------------------------------------------------------------

運動量移流の詳細定式（FCCD 面ジェット $\mathcal{J}_f(u)=(u_f,u'_f,u''_f)$ と
HFE 片側状態の境界閉包）は §\ref{sec:fccd_advection} に集約済み．
純 FCCD DNS 構成では，
方向別 HFE 状態 $(u_f^+, u_f^-)$ 上の風上選択を採用する：
\begin{equation}
  u_f^\text{upwind}
    = \begin{cases}
        u_f^- & (u_f \ge 0),\\
        u_f^+ & (u_f < 0),
      \end{cases}
  \qquad
  u_f^{+/-} = u_f \pm \tfrac{h}{2}u'_f + \tfrac{h^2}{8}u''_f
  \label{eq:pure_fccd_hfe_upwind}
\end{equation}
HFE は再構成層であり，高次ステンシルが不連続な相データを
正しい片側状態なしにサンプリングしないことを保証する．

% -------------------------------------------------------------------------------
\subsubsection{第3層：粘性項の応力発散 + DC 分離}
\label{sec:pure_fccd_phase3}
% -------------------------------------------------------------------------------

粘性項は応力テンソル発散形 $\mathcal{V}(\bu) = \bnabla\cdot(2\mu\,D(\bu))$ で保持する
（$\mu\bnabla^2\bu$ への畳み込み禁止：§\ref{sec:viscous_mu_lap_forbidden}）．
3 層構造（相内 / 界面帯 / 陰的剛性）の詳細定式は §\ref{sec:viscous_3layer} に，
DC 分離による高次残差と低次補正の階層は §\ref{sec:time_viscous} に集約済み．
純 FCCD DNS では陰的剛性層で
\begin{equation}
  A_L\,\delta\mathbf{u} = r_H,
  \qquad r_H = b_\nu - A_H\,\mathbf{u}^{(m)}
  \label{eq:pure_fccd_viscous_dc}
\end{equation}
を解く．ここで $b_\nu$ は BDF2 時間離散化・陽的項・境界条件を集めた
粘性 Helmholtz 系の右辺である．この分離は
圧力側の戦略（§\ref{sec:dc_pure_fccd_outer_inner}）と\textbf{双対}を成す
（安定な内殻解 + 高次外殻物理）．

% -------------------------------------------------------------------------------
\subsubsection{第4層：分相 FCCD PPE + GFM の中核構成}
\label{sec:pure_fccd_phase4}
% -------------------------------------------------------------------------------

圧力投影は\textbf{純 FCCD DNS の設計上の中核}である（§\ref{sec:pure_fccd_split_ppe}）．
一括 smoothed-Heaviside PPE や有限体積法フラックス釣合ではなく，
§\ref{sec:split_ppe_framework} 式~\eqref{eq:split_ppe} の分相 PPE
$\nabla^2 \chi_k = \rho_k\, q$（$k\in\{\ell,g\}$；定数密度）を
各相内部で独立に解く．
等価な変動係数形式
\begin{equation}
  \bnabla\cdot\!\left(\frac{1}{\rho_k}\bnabla \chi_k\right) = q
  \quad (k\in\{\ell,g\})
  \label{eq:pure_fccd_split_ppe_core}
\end{equation}
は，定数 $\rho_k$ により式~\eqref{eq:split_ppe} と同一（積の微分則の第 2 項は消滅）．
相は $\Gamma$ 上で GFM ジャンプ条件
（式~\eqref{eq:split_ppe_pressure_jump}--\eqref{eq:split_ppe_flux_jump}）
\begin{equation}
  [\chi]_\Gamma = J_p^\chi,
  \qquad
  \left[\frac{1}{\rho}\partial_n \chi\right]_\Gamma = J_n
  \label{eq:pure_fccd_gfm_jump}
\end{equation}
で縫合される（全圧力評価では $J_p^\chi=j_{gl}=-\sigma\kappa_{lg}$，
IPC 増分評価では $J_p^\chi=\delta j_{gl}$；非相変化非圧縮では $J_n=0$）．
FCCD が面基準点の勾配・発散演算子を供給し，
GFM がゴースト圧力ジェット $(p^*,p'^*,p''^*)$ を提供することで
コンパクトステンシルが各相側で単相のまま維持される．
離散方程式は，$\Gamma$ から離れた位置で単相高次構造を保ち，
界面不連続を物理的に導出されたジャンプ閉式のみを通して受け取る．

% -------------------------------------------------------------------------------
\subsubsection{CFD 上の位置付け：発展的 DNS / シャープ界面極限}
\label{sec:pure_fccd_positioning}
% -------------------------------------------------------------------------------

純 FCCD DNS は\textbf{発展的 DNS / シャープ界面極限}に位置づけられる：

\begin{itemize}[noitemsep]
  \item \textbf{適する現象}：マイクロ流体，インクジェット破断，
        気泡核生成，毛管波，合体・分裂開始など
        表面張力支配現象．
  \item \textbf{設計意図}：圧力・表面張力・HFE・界面ジャンプ評価を
        共通の高次面基準点言語で揃え，
        寄生流れ評価~\cite{Popinet2009,Renardy2002}で必要となる
        離散力の位置一致を保つ．
  \item \textbf{主な代償}：離散化と界面条件の複雑度，
        界面行列条件の構成難度，
        楕円型問題の条件付けと界面条件の複雑化．
  \item \textbf{対象外}：有限体積法の頑健性と正確な質量保存が第一要件の
        広範な産業スループット．
\end{itemize}

純 FCCD は二相標準経路（分相 PPE + HFE + DC）や
単相・低密度比の縮約/対照経路と競合せず\textbf{補完}する．
標準経路は汎用検証・標準計算を担い，純 FCCD は極限レジームを見据えた発展的構成である．

% -------------------------------------------------------------------------------
\subsubsection{検証章との接続}
\label{sec:pure_fccd_gates}
% -------------------------------------------------------------------------------

純 FCCD DNS は本稿の検証済み標準経路そのものではなく，
分相 PPE，HFE，FCCD 面ジェット，Balanced--Force 条件を
同じ面基準点へ集約する発展的構成である．
したがって，本節で主張できる範囲は
「どの検証済み構成要素に依存するか」を明示することに限られる．

\begin{center}
\small
\begin{tabular}{@{}p{0.26\linewidth}p{0.43\linewidth}p{0.24\linewidth}@{}}
\hline
発展的構成の要素 & 対応する検証済み構成要素 & 本稿での扱い \\
\hline
分相圧力投影 & 分相 PPE + DC + HFE & 有限時間縮約系で検証 \\
界面越し延長 & HFE 単体精度 & 1 次元 6 次，2 次元全テンソル評価で低次化なし \\
面基準点整合 & FCCD 面値・面勾配 & 面評価は 4 次律速を明示 \\
力の釣合い & Balanced--Force 静的液滴 & 静的・縮約条件で検証 \\
\hline
\end{tabular}
\end{center}

表に含まれない振動液滴固有モードや極低 We の寄生流れ評価は，
本稿では検証済み結果として扱わない．
それらは本発展的構成の適用候補であり，結論では未検証範囲として位置付ける．

% -------------------------------------------------------------------------------
\subsubsection{標準経路との関係}
\label{sec:pure_fccd_roadmap}
% -------------------------------------------------------------------------------

二相標準経路は，§9 の圧力ジャンプ閉包
（分相 PPE + ジャンプ補正付き面勾配 + HFE + DC + 大域ゲージ）を用いる構成である．
一括 smoothed-Heaviside PPE は単相・低密度比縮約または制約診断の対照経路として扱う．
純 FCCD DNS は標準経路を置き換えるものではなく，
すべての界面演算を FCCD 面ジェットへ寄せることで
シャープ界面極限を狙う発展的構成である．
本稿の検証結果が直接支える範囲は，
分相 PPE + HFE と Balanced--Force 整合を組み合わせる部分であり，
完全な第1--4層の同時使用は本稿の検証範囲外である．

% ═══════════════════════════════════════════════════════════════════════════════
